\chapter{INTRODUÇÃO}

O avanço da tecnologia digital transformou a forma como as pessoas acessam informação e se comunicam. Antes o acesso à Internet dependia principalmente de computadores de mesa e hoje em dia os \textit{smartphones} se consolidaram como o principal dispositivo de conexão. Sua portabilidade e facilidade de uso permitem acesso constante a serviços de comunicação, redes sociais, entretenimento e ferramentas acadêmicas e profissionais.

Apesar dos benefícios, o uso frequente e prolongado desses dispositivos tem sido associado a efeitos negativos. O fácil acesso ao celular ao longo de todo o dia pode levar ao uso excessivo, muitas vezes sem que o usuário perceba o tempo desperdiçado. Mesmo que esse comportamento não seja caracterizado como um vício, ainda assim esse padrão pode gerar impactos relevantes na saúde e na qualidade de vida.

Estudos indicam associação entre o uso excessivo de \textit{smartphones} e problemas como distúrbios do sono, dificuldades de atenção, aumento de ansiedade, sintomas depressivos e sedentarismo \cite{chen2023, karila2025}. Também são observados prejuízos nas relações interpessoais e no desempenho acadêmico ou profissional \cite{christensen2016}. Esses efeitos atingem diferentes áreas da vida cotidiana, incluindo aspectos físicos, cognitivos, emocionais e sociais.

Outro fator relevante é o uso de mecanismos de engajamento em aplicações digitais. Muitos aplicativos utilizam estratégias de interface conhecidas como \textit{design} persuasivo \cite{chen2023}, que incluem notificações frequentes, recompensas variáveis e personalização de conteúdo. Essas estratégias buscam manter o usuário ativo por mais tempo e incentivá-lo a acessar o aplicativo com frequência, o que pode contribuir para padrões de uso prolongados.

Diante desse cenário, fabricantes passaram a incorporar ferramentas nativas de monitoramento e controle de tempo de tela. Esses recursos permitem visualizar o tempo de uso e estabelecer limites para aplicativos específicos. No entanto, tais mecanismos dependem da iniciativa do usuário e podem ser facilmente desativados, o que limita seus resultados.

Considerando esses aspectos, surge a necessidade de investigar soluções que auxiliem na redução do tempo de uso de forma mais estruturada. Este trabalho parte da hipótese de que uma aplicação pode contribuir para o monitoramento e a redução do tempo de tela ao combinar acompanhamento automático com estratégias de engajamento voltadas à mudança de comportamento.

A abordagem adotada consiste no desenvolvimento de um aplicativo com foco na redução coletiva do tempo de uso. A solução utiliza elementos de gamificação, nos quais usuários participam de desafios semanais em grupo com o objetivo de manter o menor tempo de tela ao longo do período. A proposta utiliza mecanismos de comparação social e dinâmica coletiva como forma de incentivar os participantes a reduzirem e controlarem melhor o próprio tempo de uso do \textit{smartphone}.

\section{Fundamentação e direcionamento da proposta}

Diante dos impactos negativos já vistos na literatura sobre o uso excessivo de \textit{smartphones}, que afetam a saúde física, mental e social dos usuários, este trabalho busca responder à seguinte pergunta:

\begin{quote}
\textbf{Uma solução tecnológica pode ser eficaz no monitoramento e na redução do tempo de tela de \textit{smartphones}, visando mitigar os impactos negativos do uso excessivo na saúde e qualidade de vida dos usuários?}
\end{quote}

Para responder a essa pergunta, este trabalho estabelece como \textbf{objetivo geral} o desenvolvimento do aplicativo Desconecta, uma solução tecnológica que auxilie na redução do tempo de tela de forma eficaz, promovendo o bem-estar digital dos usuários.

Os objetivos específicos são:

\begin{enumerate}[label=\Roman*.]
    \item Analisar criticamente a literatura sobre os impactos do uso excessivo de \textit{smartphones} nas esferas física, mental, cognitiva e social;
    \item Identificar as principais limitações das ferramentas nativas de controle de tempo de tela;
    \item Definir funcionalidades e requisitos para um aplicativo que utilize estratégias gamificadas como estímulo à mudança de comportamento;
    \item Avaliar o potencial da solução proposta na promoção de hábitos digitais mais saudáveis.
\end{enumerate}

A abordagem adotada consiste no desenvolvimento de um aplicativo com foco na redução coletiva do tempo de tela, utilizando elementos de gamificação. Usuários participarão de desafios semanais em grupo, competindo entre si para manter o menor tempo de uso, em uma dinâmica inspirada em comunidades \textit{fitness} como os ``\textit{gym rats}''\footnote{Gym Rats. Disponível em: \url{https://www.gymrats.com.br/}. Acesso em: 17 fev. 2026.}. Essa competição amigável busca incentivar o senso de pertencimento e engajamento coletivo como forma de estimular uma mudança sustentável no comportamento digital.

Este trabalho adota a metodologia de pesquisa aplicada com abordagem mista, estruturada em três etapas. Na primeira etapa, foi criado um grupo de usuários no aplicativo \textit{WhatsApp}\footnote{WhatsApp. Disponível em: \url{https://www.whatsapp.com/}. Acesso em: 17 fev. 2026.} para medir o tempo de uso do celular por meio de capturas de tela dos relatórios de uso fornecidos pelo sistema. Nessa fase, também foram levantados requisitos a partir das necessidades e sugestões dos participantes. Na segunda etapa, foi realizado o projeto e a implementação da solução proposta, considerando os dados coletados e os requisitos identificados anteriormente. Na terceira etapa, o aplicativo foi disponibilizado ao mesmo grupo de usuários para testes, sendo analisadas as funcionalidades implementadas, a usabilidade da solução desenvolvida e os \textit{feedbacks} obtidos, com o objetivo de avaliar a proposta na prática.

[resultados obtidos]

\section{Organização do texto}

O restante deste trabalho está organizado da seguinte forma. No Capítulo 2, é apresentada a revisão da literatura sobre os impactos do uso excessivo de \textit{smartphones} e as limitações das ferramentas atuais de controle. No Capítulo 3, são descritos os procedimentos da pesquisa com usuários, incluindo a coleta de dados comportamentais e o levantamento de requisitos do aplicativo. No Capítulo 4, detalha-se o projeto e a implementação da solução proposta, abordando a arquitetura do sistema, as tecnologias utilizadas e as funcionalidades desenvolvidas. Por fim, no Capítulo 5, são apresentados os resultados dos testes realizados com o grupo de usuários, bem como a avaliação das funcionalidades implementadas, a análise dos \textit{feedbacks} e os demais resultados obtidos.

\bigskip
